\mypara{Visual concepts associated with human preference}
We apply frozen FG-CLIP2 dense image features and the full supplied concept bank to human-ranked alternatives within each dataset. For every concept, we average preferred-minus-rejected maximum patch similarities within annotation events, then within normalized prompts, then equally across prompts. ImageReward uses all strict ranked pairs; the linked four-candidate Pick-a-Pic release uses the selected image against each non-selected alternative. We freeze concepts and directions on discovery data and evaluate held-out prompts, preserving prompt and image dependencies. We report paired concordance and component-bootstrap uncertainty, with best-versus-rest and generator/patch-count sensitivities. These are associations with recorded choices, not calibrated prevalence estimates or causal explanations. This engineering pilot retained 75 held-out prompts in 75 dependence components for imagereward. Table entries report observed effects and per-concept intervals; visual interpretations await human review.
The first discovery-ranked terms were lip-syncing, lip-smacking, stare. Of 20 frozen candidates, 20 retained their discovery direction in held-out effects; this descriptive count is not a multiplicity-adjusted significance claim. The equal signed mean of the selected concepts achieved held-out paired concordance 0.597 (95\% component-bootstrap interval [0.490, 0.697]). These pilot findings remain provisional: visual concept presence is unverified, the supplied text-bank encoding provenance is unresolved, and the sample is small.
